\section{Discussion and limitations}
\label{sec:discussion}

\paragraph{What the results mean for benchmark users.}
A one- or two-point NAVSIM gain can come from plan content that a deployed, replanning system never
executes: our benign perturbations alone cost policies \num{TailDropMinNavsimVone}--\num{TailDropMaxNavsimVone}
PDMS points.
Conversely, a planner can lead NAVSIM and trail simpler baselines in closed loop, as DiffusionDrive
does.
We do not claim that NAVSIM rankings are wrong in general; we claim that they reward something closed
loop does not, and that this is measurable.
Two practical consequences follow: plan quality beyond the replanning horizon should be scored
separately or down-weighted, and evaluators should be checked for specificity as well as sensitivity.

\paragraph{The reference is a simulator.}
BridgeSim renders camera input in MetaDrive, so camera-based policies face a rendering gap that the
blind Ego MLP does not.
This is the most plausible explanation for the Ego MLP scoring above TransFuser and DiffusionDrive in
closed loop, and we do not treat that comparison as evidence against NAVSIM.
Three of our findings do not depend on it.
The specificity test compares each policy with itself under edits that the controller provably never
executes, so a rendering gap affects both sides equally.
The degradation orders are validated per policy, with the same rendering on every level.
And LTF, which uses cameras only, is the best closed-loop policy, so rendering does not simply penalise
camera input.
Our TransFuser reproduction lies below the published BridgeSim EPDMS for all three seeds, while its DS
is close (\num{DsTFNavhard} vs.\ \num{DsTFPublished}); we report this rather than tuning the protocol.

\paragraph{One reference.}
Every ground-truth statement rests on one closed-loop simulator with log-replay traffic.
Its rankings are stable across two disjoint scene sets ($\tau={}$\num{TauNavhardHeldout}) and its
per-scene scores are reproducible (Spearman \num{RetestMin}--\num{RetestMax}), but a second reference
(a reconstruction-based simulator or real-world tests) would make the ground truth stronger.
HUGSIM~\cite{zhou2024hugsim} needs compiled extensions our cluster does not allow, and no openly
available world model is both multi-camera and action-conditioned for NAVSIM policies at the memory of
our GPUs; both are left to future work.

\paragraph{Specificity depends on the controller.}
The benign perturbations are benign for a controller that replans every $0.5$\,s and tracks the first
seconds of a plan.
For a system that executes longer horizons, the tail of a plan matters; the test should then be
re-derived from that system's controller.

\paragraph{Choices made after seeing data.}
We disclose three.
The closed-loop score is DS rather than BridgeSim's EPDMS, chosen after observing that EPDMS rewards a
policy for driving less; \cref{tab:robustness} reports EPDMS.
The disjoint-scene design was adopted after the half-split analysis lost power
(\cref{sec:results:robust}); its measurement scenes were fixed before they were simulated, and the
split analyses are reported.
One measurement scene that the simulator could not run for any policy was excluded for all.
The measurement scenes are ordinary navtest scenes, easier than NavHard, so absolute closed-loop
scores are higher there; all comparisons are within a scene set.

\section{Conclusion}
We tested driving evaluators against rankings that are known by construction and validated in closed
loop on separate scenes.
Only latency and lateral drift yield such rankings; perception corruptions do not, even in closed loop.
On the validated failures, the most sensitive evaluator is open-loop displacement, yet it and both
NAVSIM versions penalise plan changes that a replanning vehicle never executes, which closed loop
ignores.
On real policies, NAVSIM and closed loop disagree on held-out scenes far beyond simulator noise.
Sensitivity to known failures is not enough to validate an evaluator; specificity has to be tested too,
and our harness makes both tests cheap to run on any new benchmark.
